\section{Challenges}
\label{sec:intro_challenges}

The challenges listed here surfaced repeatedly during the literature review, and they are expected to shape the work over the next six months regardless of which specific design choices are ultimately made. They appear here, at an early stage, rather than waiting to encounter them during implementation:

\begin{enumerate}
	\item \textbf{Inferring dependencies without ground truth.} Any prerequisite structure the agent builds on its own is inferred by an LLM rather than annotated by an expert, and this is expected to be the single riskiest part of whatever gets built. Before relying on such a structure for any evaluation, the plan is to run a small pilot study against PentestEval's expert-annotated dependencies to get a sense of how accurate the inference actually is.
	\item \textbf{Sandbox results may not transfer to the real world.} \citet{fang2024hackwebsites} found only 1 exploitable XSS in 50 real-world candidate sites (2\%), against 73.3\% in a matched sandbox. This gap is worth keeping in mind throughout, to avoid presenting sandbox numbers as if they were real-world numbers.
	\item \textbf{Some benchmarks are small.} PrediQL's GraphQL suite, for instance, covers only 6 APIs --- real and standardised, but far smaller than benchmarks like CVE-Bench. Any conclusion drawn from it needs to be stated as exploratory rather than definitive.
	\item \textbf{Cost is a moving target.} Whatever cost-per-exploit figures are eventually reported, they are tied to a specific model and pricing at a specific moment, and that needs to be stated explicitly rather than presenting cost as a fixed property of the approach.
	\item \textbf{Reused strategies could backfire.} A strategy that worked against one version of a piece of software could plausibly be actively harmful against a different version. The reviewed papers do not appear to address this risk, which is one reason the design treats it as worth attention from the start.
	\item \textbf{Sensitivity to configuration choices.} If the approach depends on tunable parameters, checking how sensitive results are to those choices matters, rather than reporting numbers from a single, possibly fortunate, configuration.
	\item \textbf{Honesty about what generalises.} If testing spans more than one attack surface, it's important to distinguish between ``a shared design'' and ``one literally unmodified system'', since surface-specific components are expected to be necessary either way.
\end{enumerate}
